\documentclass[10pt,a4paper]{amsart}
\usepackage[margin=19mm]{geometry}
\usepackage{amsmath,amssymb,mathtools}
\usepackage[scr=rsfs,cal=euler]{mathalfa}
\usepackage{xfrac}
\usepackage[unicode,hidelinks,bookmarks=false]{hyperref}
\newcommand{\ie}{i.e.\ }
\newcommand{\cf}{cf.\ }
\newcommand{\resp}{respectively\ }
\newcommand{\Subsection}{\subsection}
\providecommand{\nobreakdash}{\nobreak\hbox{-}}
\setlength{\emergencystretch}{2em}
\multlinegap0pt
\usepackage{bm}
\usepackage{bbm}
\usepackage{dsfont}
\usepackage{xspace}
\usepackage{cancel}
\usepackage{mathtools}
\usepackage{amsthm}
\makeatletter
\@ifundefined{thm}{\newtheorem{thm}{Thm}}{}
\@ifundefined{lem}{\newtheorem{lem}[thm]{Lemma}}{}
\makeatother
\makeatletter
\newcommand{\HH}{\mathscr H}
\def\bz{\beta}
\def\ez{\epsilon}
\expandafter\ifx\csname note\endcsname\relax
\newenvironment{note}{\begin{quote}\sf}{\end{quote}}
\fi
\def\sz{\sigma}
\makeatother
\title[Sharp stability of Alexandrov's theorem for $C^1$ domains in]{Sharp stability of Alexandrov's theorem for $C^1$ domains in the small-excess regime}
\author{Alessio Figalli, Yi Ru-Ya Zhang}
\date{Source: arXiv:2606.13335v1 (CC BY 4.0)}
\begin{document}
\maketitle

\noindent\emph{Source and licence.} Sharp stability of Alexandrov's theorem for $C^1$ domains in the small-excess regime. Version arXiv:2606.13335v1 (CC BY 4.0), distributed under the Creative Commons Attribution 4.0 International licence, as recorded in the arXiv licence metadata for this exact version and shown on the abstract page https://arxiv.org/abs/2606.13335v1. Full rights notice: licenses/arxiv-2606.13335-CC-BY-4.0.txt.\par\smallskip

\noindent\textit{Editorial scope.} Counted unit: the Lem whose source label is \texttt{strong l2} in the article (source0.tex lines 754--777, source label \texttt{strong l2}), delivered together with the article's own complete original proof of it (source0.tex lines 779--1276, one \texttt{proof} environment of 498 lines). No mathematics is written, repaired or re-derived: statement and proof are the article's own text line for line. Required context retained uncounted: none. Every definition, display and label of those lines is retained unchanged. Omitted from this package and not redistributed: 1--753, 778--778, 1277--3773. Nothing inside a retained excerpt is omitted or altered: every retained body is delivered exactly as the article's lines are written, so no omission receipt of any kind accompanies this package. Citations: the retained excerpts carry no citation command at all, so no bibliography entry of the article is transcribed and no reference is redistributed. Reference adaptation: none was needed and none was made; every reference inside the delivered unit resolves inside it, so no unresolved reference or citation is produced. Printed slips kept literal: no printed word, symbol, hypothesis or number of the counted statement or of its proof was altered. Layout and class: the article's own class is \texttt{amsart} with packages including a4wide, color, eucal, enumerate, mathrsfs, ulem, amsmath, amssymb, epsfig, amsthm, inputenc, psfrag, hyperref, cleveref; the wrapper is the lane's pinned \texttt{amsart} editorial wrapper at 10pt on A4, which restates the theorem-like environments the retained text needs and adds the article's own macro definitions that the retained text needs (\texttt{\textbackslash HH}, \texttt{\textbackslash bz}, \texttt{\textbackslash ez}, \texttt{\textbackslash sz}), with extra wrapper packages (bm, bbm, dsfont, xspace, cancel, mathtools). The delivered numbering is the unit's own. Assets: no retained span contains a figure environment, a graphics-inclusion command, a PDF-page inclusion, a table or a TikZ picture; no image file is copied into this package, the package declares no asset dependency and no image file is redistributed with it. Licence: version arXiv:2606.13335v1 of this article is distributed under the Creative Commons Attribution 4.0 International licence, as recorded in the arXiv licence metadata for this exact version and shown on the abstract page https://arxiv.org/abs/2606.13335v1. A full rights notice is delivered at licenses/arxiv-2606.13335-CC-BY-4.0.txt. Characters: every retained excerpt is pure ASCII, so the delivered engine, XeLaTeX with its default Latin Modern text font, reports no missing character.
\begin{lem}\label{strong l2}
Let $\xi \in {BV(\mathbb S^{n-1})}$ satisfy, for some  $0<\sz=\sz(n)\ll1$,
\begin{equation}\label{small average}
    \left|\frac{1}{ n|B|}\int_{\mathbb S^{n-1}}\xi\,  d\mathscr H^{n-1}\right|\le \sz \left(\int_{\mathbb S^{n-1}} \frac{|\xi|^2}{\sqrt{1+\ez^2|\xi|^2}} \,d\mathscr H^{n-1}\right)^{\frac 1 2}.
\end{equation}
 Then,  for any $0<\ez<\frac 1 2,$
$$   \int_{\mathbb S^{n-1}} \frac{|\xi|^2}{\sqrt{1+\ez^2|\xi|^2}} \,d\mathscr H^{n-1}\le C(n)\left( \int_{\mathbb S^{n-1}} \frac{|\nabla_\tau  \xi|^2}{\sqrt{1+\ez^2|\nabla_\tau \xi|^2}} \,
 d\mathscr H^{n-1}+\ez^{-1} |D^s_\tau \xi|(\mathbb S^{n-1})\right).$$
In addition, assume that there exists a sequence $\xi_i\in BV(\mathbb S^{n-1})$ satisfying, for every $i$,
\begin{equation}\label{small average seq}
    \left|\frac{1}{ n|B|}\int_{\mathbb S^{n-1}}\xi_i\,  d\mathscr H^{n-1}\right|\le \sz \left(\int_{\mathbb S^{n-1}} \frac{|\xi_i|^2}{\sqrt{1+\ez_i^2|\xi_i|^2}} \,d\mathscr H^{n-1}\right)^{\frac 1 2},
\end{equation}
and
$$\int_{\mathbb S^{n-1}} \frac{|\nabla_\tau  \xi_i|^2}{\sqrt{1+\ez_i^2|\nabla_\tau \xi_i|^2}} \,
 d\mathscr H^{n-1}+ \epsilon_i^{-1}|D^s_\tau \xi_i|(\mathbb S^{n-1})\le M$$
for some $\ez_i\to 0$ and $M>0$.
Then, up to passing to a subsequence,
$\xi_i \to \xi$   strongly in $L^p (\mathbb S^{n-1})$ for some $1<p=p(n)<2$ (below the BV-Sobolev compactness exponent on $S^{n-1}$),
and
\begin{equation}
    \label{eq:xi_i to xi}
\lim_{i\to \infty}\int_{\mathbb S^{n-1}} \frac{|\xi_i|^2}{\sqrt{1+\ez_i^2|\xi_i|^2}} \,d\mathscr H^{n-1}=  \int_{\mathbb S^{n-1}} {|\xi|^2}  \,d\mathscr H^{n-1}.
\end{equation}
\end{lem}

\begin{proof}
\noindent{\bf Step 1: Energy inequalities.}
Given $\xi\in BV(\mathbb S^{n-1})$,  we set for simplicity
$$
\mathcal A_\epsilon(\xi):=\int_{\mathbb S^{n-1}} \frac{|\xi|^2}{\sqrt{1+\epsilon^2|\xi|^2}}\,d\mathscr H^{n-1},
$$
$$
\mathcal E_\epsilon(\xi):=\int_{\mathbb S^{n-1}} \frac{|\nabla_\tau \xi|^2}{\sqrt{1+\epsilon^2|\nabla_\tau \xi|^2}}\,d\mathscr H^{n-1}
+\ez^{-1}|D^s_\tau \xi|(\mathbb S^{n-1}).
$$
We first prove the estimate for smooth functions. The passage to $BV$ follows by strict approximation on the sphere. Indeed, after subtracting constants, one may choose
$\xi_h\in C^\infty(\mathbb S^{n-1})$ with
$\int\xi_h=\int\xi$, $\xi_h\to\xi$ in $L^1$, and
$D_\tau\xi_h\stackrel{*}{\rightharpoonup}D_\tau\xi$ with
$|D_\tau\xi_h|(\mathbb S^{n-1})\to |D_\tau\xi|(\mathbb S^{n-1})$. Since the
recession function of $p\mapsto |p|^2(1+\epsilon^2|p|^2)^{-1/2}$ is
$\epsilon^{-1}|p|$, Reshetnyak's theorem gives
$\mathcal E_\epsilon(\xi_h) \to \mathcal E_\epsilon(\xi)$ and
$\mathcal A_\epsilon(\xi_h)\to\mathcal A_\epsilon(\xi)$. We may therefore work
in the smooth setup.

We first consider the case where $\sz=0$ in \eqref{small average}.

\medskip
\noindent{\bf Case 1}: $\|\ez \xi\|_{L^\infty(\mathbb S^{n-1})}\le 1$. Then
\begin{equation}\label{comparable}
 \frac 1 2 \int_{\mathbb S^{n-1}}  {|\xi|^2}  \,d\mathscr H^{n-1}\le \int_{\mathbb S^{n-1}} \frac{|\xi|^2}{\sqrt{1+\ez^2|\xi|^2}} \,d\mathscr H^{n-1}\le \int_{\mathbb S^{n-1}}  {|\xi|^2}  \,d\mathscr H^{n-1}.
\end{equation}
We claim that, for any $0<\bz<1$,
\begin{equation}\label{youngs}
|\xi||\nabla_\tau \xi|\le C(\bz) \frac{|\nabla_\tau  \xi|^2}{\sqrt{1+\ez^2|\nabla_\tau \xi|^2}} + \bz |\xi|^2.
\end{equation}
Indeed, if $\ez|\nabla_\tau \xi|\le 1$, then by Young's inequality,
$$|\xi||\nabla_\tau \xi|\le C(\beta)  |\nabla_\tau \xi|^2 +\bz |\xi|^2\le \sqrt{2} C(\bz) \frac{|\nabla_\tau  \xi|^2}{\sqrt{1+\ez^2|\nabla_\tau \xi|^2}} + \bz |\xi|^2. $$
When $\ez |\nabla_\tau \xi|\ge 1$, it follows from $\ez|\xi|\le 1$ that
$$|\xi||\nabla_\tau \xi|\le \ez^{-1} |\nabla_\tau \xi|\le \sqrt{2}  \frac{|\nabla_\tau  \xi|^2}{\sqrt{1+\ez^2|\nabla_\tau \xi|^2}}. $$
Thus, we conclude the validity of \eqref{youngs}.

Now, by applying a $(1,1)$-Poincar\'e inequality and \eqref{youngs} (recall that $\sz=0$ in \eqref{small average}), we get
\begin{align*}
\int_{\mathbb S^{n-1}} |\xi|^2 \,d\mathscr H^{n-1}& \leq  C(n) \int_{\mathbb S^{n-1}} |\xi \nabla_\tau \xi| \,d\mathscr H^{n-1}\\
& \leq  C(n,\bz) \int_{\mathbb S^{n-1}}  \frac{|\nabla_\tau  \xi|^2}{\sqrt{1+\ez^2|\nabla_\tau \xi|^2}}   + C(n)\bz \int_{\mathbb S^{n-1}}   |\xi|^2 \,d\mathscr H^{n-1}.
\end{align*}
Then by taking $\bz>0$ small enough, we obtain
\begin{equation}\label{xi2 L2}
 \int_{\mathbb S^{n-1}} |\xi|^2 \,d\mathscr H^{n-1} \le  C(n)  \int_{\mathbb S^{n-1}} \frac{|\nabla_\tau  \xi|^2}{\sqrt{1+\ez^2|\nabla_\tau \xi|^2}} \,d\mathscr H^{n-1},
\end{equation}
and also
\begin{equation}\label{xi2 W11}
    \int_{\mathbb S^{n-1}} |\xi \nabla_\tau \xi| \,d\mathscr H^{n-1}\le  C(n)  \int_{\mathbb S^{n-1}} \frac{|\nabla_\tau  \xi|^2}{\sqrt{1+\ez^2|\nabla_\tau \xi|^2}} \,d\mathscr H^{n-1}.
\end{equation}
Recalling \eqref{comparable}, this proves the first part of the lemma under the assumption of Case 1.

\medskip
\noindent{\bf Case 2}: General case.
Define the truncation map
$$
T_\epsilon(t):=\max\left\{-\frac1\epsilon,\min\Bigl\{t,\frac1\epsilon\Bigr\}\right\}
$$
and consider the function
$$
F(c):=\int_{\mathbb S^{n-1}} T_\epsilon(\xi-c)\,d\mathscr H^{n-1}.
$$
Since $T_\epsilon$ is continuous and nondecreasing, $F$ is continuous and nonincreasing.
Moreover,
$$
\lim_{c\to-\infty}F(c)=\frac{\mathscr H^{n-1}(\mathbb S^{n-1})}{\epsilon}>0,
\qquad
\lim_{c\to+\infty}F(c)=-\frac{\mathscr H^{n-1}(\mathbb S^{n-1})}{\epsilon}<0,
$$
hence the intermediate value theorem yields the existence of $c_\epsilon\in\mathbb R$ for which 
$$
\int_{\mathbb S^{n-1}} T_\epsilon(\xi-c_\epsilon)\,d\mathscr H^{n-1}=0.
$$
Set
$$
\tilde\xi:=\xi-c_\epsilon,\qquad
\psi:=T_\epsilon(\tilde\xi),\qquad
\eta:=\tilde\xi-\psi
=
\Bigl(\tilde\xi-\frac 1 \ez \Bigr)_+-\Bigl(\tilde\xi+\frac 1 \ez \Bigr)_-.
$$
Then
$$
\tilde\xi=\psi+\eta,\qquad \int_{\mathbb S^{n-1}}\psi\,d\mathscr H^{n-1}=0,
\qquad \|\epsilon\psi\|_{L^\infty(\mathbb S^{n-1})}\le 1.
$$
Also,
$$
\nabla_\tau\psi=\left\{
\begin{array}{ll}
\nabla_\tau\xi& \text{$\mathscr H^{n-1}$-a.e. on $\mathbb S^{n-1}\cap \{\ez|\tilde\xi|\leq  1\}$},\\
0& \text{$\mathscr H^{n-1}$-a.e. on $\mathbb S^{n-1}\cap \{\ez|\tilde\xi|> 1\}$},
\end{array}\right.
$$
and similarly
$$
\nabla_\tau\eta=\left\{
\begin{array}{ll}
0& \text{$\mathscr H^{n-1}$-a.e. on $\mathbb S^{n-1}\cap \{\ez|\tilde\xi|\leq  1\}$},\\
\nabla_\tau\xi& \text{$\mathscr H^{n-1}$-a.e. on $\mathbb S^{n-1}\cap \{\ez|\tilde\xi|> 1\}$}.
\end{array}\right.
$$
Since
$$
\left|\frac{d}{dt}\frac{|t|^2}{\sqrt{1+\epsilon^2|t|^2}}\right|
\le C\,\epsilon^{-1}\qquad\forall\,t\in\mathbb R,
$$
we have
$$
\mathcal A_\epsilon(\xi)\le \mathcal A_\epsilon(\tilde\xi)+C\epsilon^{-1}|c_\epsilon|.
$$
On the other hand, since $\int_{\mathbb S^{n-1}}\xi\,d\mathscr H^{n-1}=0$ and
$\int_{\mathbb S^{n-1}}\psi\,d\mathscr H^{n-1}=0$, we get
$$
-c_\epsilon\,\mathscr H^{n-1}(\mathbb S^{n-1})
=\int_{\mathbb S^{n-1}}\tilde\xi\,d\mathscr H^{n-1}
=\int_{\mathbb S^{n-1}}\eta\,d\mathscr H^{n-1},
$$
hence
$$
|c_\epsilon|\le C(n)\int_{\mathbb S^{n-1}}|\eta|\,d\mathscr H^{n-1}.
$$
Therefore
$$
\mathcal A_\epsilon(\xi)\le \mathcal A_\epsilon(\tilde\xi)+C(n)\epsilon^{-1}\int_{\mathbb S^{n-1}}|\eta|\,d\mathscr H^{n-1}.
$$
We claim that, pointwise,
\begin{equation}\label{pointwise tilde xi}
\frac{|\tilde\xi|^2}{\sqrt{1+\epsilon^2|\tilde\xi|^2}}
\le
C |\psi|^2 + C\epsilon^{-1}|\eta|.
\end{equation}
Indeed, if $|\eta|\le |\psi|$, then $|\tilde\xi|\le 2|\psi|$ and $|\psi|\le \epsilon^{-1}$, so
$$
\frac{|\tilde\xi|^2}{\sqrt{1+\epsilon^2|\tilde\xi|^2}}
\le |\tilde\xi|^2 \le 4|\psi|^2.
$$
If $|\eta|>|\psi|$, then $|\tilde\xi|\le 2|\eta|$, and hence
$$
\frac{|\tilde\xi|^2}{\sqrt{1+\epsilon^2|\tilde\xi|^2}}
\le \epsilon^{-1}|\tilde\xi|\le 2\epsilon^{-1}|\eta|,
$$
which concludes the validity of \eqref{pointwise tilde xi}. 
Thus
\begin{equation}\label{two terms}
    \begin{split}
     \int_{\mathbb S^{n-1}} \frac{|\xi|^2}{\sqrt{1+\ez^2|\xi|^2}} \,d\mathscr H^{n-1}
 & \leq  C \int_{\mathbb S^{n-1}}  {|\psi|^2}  \,d\mathscr H^{n-1}+ C \ez^{-1}\int_{\mathbb S^{n-1}} |\eta|  \, \,d\mathscr H^{n-1}.
    \end{split}
\end{equation}
Also, by applying the results obtained in Case 1 to $\psi$ (recall that $\int_{\mathbb S^{n-1}}\psi\,d\mathscr H^{n-1}=0$), we get
\begin{equation}
    \begin{split}
   \int_{\mathbb S^{n-1}} |\psi|^2 \,d\mathscr H^{n-1}
  &\le  C(n)  \int_{\mathbb S^{n-1}} \frac{|\nabla_\tau  \psi|^2}{\sqrt{1+\ez^2|\nabla_\tau \psi|^2}} \,d\mathscr H^{n-1}\\
 &=  C(n)  \int_{\mathbb S^{n-1}\cap \{\ez|\tilde\xi|\le 1\}} \frac{|\nabla_\tau  \xi|^2}{\sqrt{1+\ez^2|\nabla_\tau \xi|^2}} \,d\mathscr H^{n-1}\\
  &\le  C(n)  \int_{\mathbb S^{n-1} } \frac{|\nabla_\tau  \xi|^2}{\sqrt{1+\ez^2|\nabla_\tau \xi|^2}} \,d\mathscr H^{n-1}.  \label{first part}
    \end{split}
\end{equation}
and
\begin{equation}
    \begin{split}
\int_{\mathbb S^{n-1}} |\psi \nabla_\tau \psi| \,d\mathscr H^{n-1}& \le  C(n)  \int_{\mathbb S^{n-1}} \frac{|\nabla_\tau  \psi|^2}{\sqrt{1+\ez^2|\nabla_\tau \psi|^2}} \,d\mathscr H^{n-1}  \\
& \le  C(n)  \int_{\mathbb S^{n-1} } \frac{|\nabla_\tau  \xi|^2}{\sqrt{1+\ez^2|\nabla_\tau \xi|^2}} \,d\mathscr H^{n-1}. \label{psi W11}
    \end{split}
\end{equation}
Combining \eqref{first part} and \eqref{two terms}, we see that
it suffices to show that
\begin{equation} \label{eta controlled by xi}
\ez^{-1}\int_{\mathbb S^{n-1}} |\eta|  \, \,d\mathscr H^{n-1}\le C(n)\int_{\mathbb S^{n-1}}   \frac{|\nabla_\tau  \xi|^2}{\sqrt{1+\ez^2|\nabla_\tau \xi|^2}}  \, \,d\mathscr H^{n-1}.
\end{equation}
We now claim that $0$ is the median of $\eta$. Indeed, note that
$$
\{\eta>0\}=\{\psi= \epsilon^{-1}\} \quad \text{  and }\quad  \{\eta<0\}=\{\psi=-\epsilon^{-1}\}.
$$
Then, since $\psi\in [-\epsilon^{-1},  \epsilon^{-1}]$ and $\int_{\mathbb S^{n-1}}\psi\,d\mathscr H^{n-1}=0$,   
$$ \text{ both }  \ \mathscr H^{n-1}(\{\eta>0\}) \ \text{ and } \ \mathscr H^{n-1}(\{\eta<0\})$$ 
are no more than  $\frac 1 2\mathscr H^{n-1}(\mathbb S^{n-1})$, which proves the claim. Thus,  the $(1,1)$-Poincar\'e inequality yields
$$
\int_{\mathbb S^{n-1}} |\eta| \,d\mathscr H^{n-1}
\le C(n)\int_{\mathbb S^{n-1}} |\nabla_\tau\eta|\,d\mathscr H^{n-1}.
$$

\medskip
\noindent {\bf SubCase 2.1}:
We first assume that
\begin{equation}\label{eta small}
\int_{\mathbb S^{n-1}} |\eta|  \, \,d\mathscr H^{n-1}\le \ez^{-1}|\mathbb S^{n-1}\cap \{\ez|\tilde\xi|>1\}|.
\end{equation}
Therefore, as $|\psi|=\ez^{-1}$ on $\mathbb S^{n-1}\cap\{\ez |\tilde\xi|> 1\}$,
\begin{equation}\label{eta small2}
\ez^{-1}\int_{\mathbb S^{n-1}} |\eta|  \, \,d\mathscr H^{n-1}\leq
\int_{\mathbb S^{n-1}\cap\{\ez|\tilde\xi|>1\}}  \ez^{-2} \,d\mathscr H^{n-1}\leq \int_{\mathbb S^{n-1}}  {|\psi|^2}  \,d\mathscr H^{n-1}.
\end{equation}
Thus, \eqref{first part} directly gives \eqref{eta controlled by xi}.

\medskip
\noindent {\bf SubCase 2.2}:
When
\begin{equation}\label{eta large}
\int_{\mathbb S^{n-1}} |\eta|  \, \,d\mathscr H^{n-1}\ge \ez^{-1}|\mathbb S^{n-1}\cap\{\ez|\tilde\xi|>1\}|,
\end{equation}
we write
$$\mathcal E= \{\ez|\nabla_\tau\xi|<1\}\cap\{\ez|\tilde\xi|>1\}\cap\mathbb S^{n-1}, \quad  \mathcal F= \{\ez|\nabla_\tau\xi|\ge 1\}\cap\{\ez|\tilde\xi|>1\}\cap\mathbb S^{n-1}.$$
Then, since $\nabla_\tau\eta=0$ $\mathscr H^{n-1}$-a.e. on $\mathbb S^{n-1}\cap \{\ez|\tilde\xi|\leq  1\}$, the $(1,1)$-Poincar\'e inequality implies
\begin{equation}
    \begin{split}
\ez^{-1}\int_{\mathbb S^{n-1}} |\eta|  \, \,d\mathscr H^{n-1} & \leq  C(n) \ez^{-1} \int_{\mathbb S^{n-1} }  |\nabla_\tau \eta| \, \,d\mathscr H^{n-1}\\
& =  C(n) \ez^{-1} \int_{\mathcal E}  |\nabla_\tau \eta| \, \,d\mathscr H^{n-1}+ C(n) \ez^{-1} \int_{\mathcal F}  |\nabla_\tau \eta| \, \,d\mathscr H^{n-1}. \label{D eta}
    \end{split}
\end{equation}

\medskip
\noindent {\bf SubCase 2.2.1}:
Suppose first that
\begin{equation}\label{F large}
\int_{\mathcal F}  |\nabla_\tau \eta| \, \,d\mathscr H^{n-1}\ge \int_{\mathcal E}  |\nabla_\tau \eta| \, \,d\mathscr H^{n-1}.
\end{equation}
Then, by the definition of $\mathcal F$ and \eqref{D eta},
\begin{multline}\label{control 1}
\ez^{-1}\int_{\mathbb S^{n-1}} |\eta|  \, \,d\mathscr H^{n-1}  \leq  C(n) \ez^{-1} \int_{\mathbb S^{n-1} }  |\nabla_\tau \eta| \, \,d\mathscr H^{n-1}\\
\leq  2C(n) \ez^{-1} \int_{\mathcal F}  |\nabla_\tau \eta| \, \,d\mathscr H^{n-1}\le 2\sqrt{2} C(n)  \int_{\mathcal F} \frac{|\nabla_\tau  \xi|^2}{\sqrt{1+\ez^2|\nabla_\tau \xi|^2}} \, \,d\mathscr H^{n-1},
    \end{multline}
where we used that $\nabla_\tau\eta=\nabla_\tau  \xi$ $\mathscr H^{n-1}$-a.e. on $\mathbb S^{n-1}\cap \{\ez|\tilde\xi|>  1\}$ together with the definition of $\mathcal F$.
This gives \eqref{eta controlled by xi}.

\medskip
\noindent {\bf SubCase 2.2.2}: Now the remaining case is
\begin{equation}\label{E large}
\int_{\mathcal F}  |\nabla_\tau \eta| \, \,d\mathscr H^{n-1}\le \int_{\mathcal E }  |\nabla_\tau \eta| \, \,d\mathscr H^{n-1}.
\end{equation}
Notice that  \eqref{eta large} yields
$$|\mathbb S^{n-1}\cap\{\ez|\tilde\xi|>1\}|= \int_{\mathbb S^{n-1}\cap\{\ez|\tilde\xi|>1\}} 1\, d\mathscr H^{n-1} \le \ez \int_{\mathbb S^{n-1}} |\eta|  \, \,d\mathscr H^{n-1}.$$
Therefore, since $\tilde\xi=\psi+\eta$, the definition of $\mathcal E$ gives
\begin{equation}\label{control 2}
|\mathcal E|\le \int_{\mathcal E} \ez|\tilde\xi| \,  \,d\mathscr H^{n-1}= \int_{\mathbb S^{n-1}\cap\{\ez|\tilde\xi|>1\}} (1+\ez|\eta|) \,  \,d\mathscr H^{n-1}\le 2\ez \int_{\mathbb S^{n-1}} |\eta|  \, \,d\mathscr H^{n-1}.
\end{equation}
Thus, by \eqref{D eta}, \eqref{E large}, H\"older's inequality, the definition of $\mathcal E$, and \eqref{control 2}, we get
\begin{align*}
\ez^{-1}\int_{\mathbb S^{n-1}} |\eta|  \, \,d\mathscr H^{n-1}& \leq 2 C(n) \ez^{-1} \int_{\mathcal E}  |\nabla_\tau \eta| \, \,d\mathscr H^{n-1} =2C(n) \ez^{-1} \int_{\mathcal E}  |\nabla_\tau \xi| \, \,d\mathscr H^{n-1}\nonumber\\
& \leq 2 C(n)\ez^{-1} \left(\int_{\mathcal E}   \frac{|\nabla_\tau  \xi|^2}{\sqrt{1+\ez^2|\nabla_\tau \xi|^2}}  \, \,d\mathscr H^{n-1}\right)^{1/2} \left(\int_{\mathcal E}    {\sqrt{1+\ez^2|\nabla_\tau \xi|^2}}  \, \,d\mathscr H^{n-1}\right)^{1/2}\nonumber\\
& \leq 2\sqrt{2} C(n)\ez^{-1} |\mathcal E|^{1/2} \left(\int_{\mathcal E}   \frac{|\nabla_\tau  \xi|^2}{\sqrt{1+\ez^2|\nabla_\tau \xi|^2}}  \, \,d\mathscr H^{n-1}\right)^{1/2}\nonumber\\
& \leq 4 C(n) \ez^{-1/2}  \left(\int_{\mathbb S^{n-1}}   |\eta|  \, \,d\mathscr H^{n-1} \right)^{1/2}    \left(\int_{\mathcal E}   \frac{|\nabla_\tau  \xi|^2}{\sqrt{1+\ez^2|\nabla_\tau \xi|^2}}  \, \,d\mathscr H^{n-1}\right)^{1/2}
\end{align*}
which yields \eqref{eta controlled by xi}  and
\begin{equation}\label{eta W11}
    \ez^{-1} \int_{\mathcal E}  |\nabla_\tau \eta| \, \,d\mathscr H^{n-1} \le C(n)\int_{\mathcal E}   \frac{|\nabla_\tau  \xi|^2}{\sqrt{1+\ez^2|\nabla_\tau \xi|^2}}  \, \,d\mathscr H^{n-1}.
\end{equation}
This concludes the proof of the first part of the lemma when $\sz=0$.

\medskip
We now treat the general case in \eqref{small average}. Let
$$
m:=\frac{1}{n|B|}\int_{\mathbb S^{n-1}}\xi\,d\mathscr H^{n-1},
\qquad
\bar\xi:=\xi-m.
$$
Then $\int_{\mathbb S^{n-1}}\bar\xi\,d\mathscr H^{n-1}=0$ and $\mathcal E_\epsilon(\bar\xi)=\mathcal E_\epsilon(\xi)$.
Set
$$
g_\epsilon(t):=\min\{|t|^2,\epsilon^{-1}|t|\}.
$$
One easily checks that
$$
\frac1{\sqrt2}\,g_\epsilon(t)\le \frac{|t|^2}{\sqrt{1+\epsilon^2|t|^2}}\le g_\epsilon(t)
\qquad\forall\, t\in\mathbb R,
$$
and that
$$
g_\epsilon(a+b)\le 2g_\epsilon(a)+2g_\epsilon(b)\qquad\forall\, a,b\in\mathbb R.
$$
Hence
$$
\mathcal A_\epsilon(\xi)=\mathcal A_\epsilon(\bar\xi+m)\le C \mathcal A_\epsilon(\bar\xi)+C \mathcal A_\epsilon(m)\le C \mathcal A_\epsilon(\bar\xi)+C |m|^2.
$$
Applying the already proved case $\sz=0$ to $\bar\xi$, we get
$$
\mathcal A_\epsilon(\bar\xi)\le C(n)\mathcal E_\epsilon(\bar\xi)=C(n)\mathcal E_\epsilon(\xi).
$$
On the other hand, \eqref{small average} gives
$$
|m|^2\le \sz^2 \mathcal A_\epsilon(\xi).
$$
Choosing $\sz=\sz(n)>0$ sufficiently small, we can absorb the $|m|^2$ term into the left-hand side and conclude the first part of the lemma in full generality.


\medskip

\noindent{\bf Step 2: Energy convergence.} In this second part, we argue directly in $BV$.
By the first part, the quantities
$$
\mathcal  A_{\epsilon_i}(\xi_i)
$$
are uniformly bounded. In particular, \eqref{small average seq} yields a uniform bound on the averages
$$
m_i:=\frac{1}{n|B|}\int_{\mathbb S^{n-1}}\xi_i\,d\mathscr H^{n-1}.
$$
Also, since $\epsilon_i<1$ for $i$ large,
$$
|a|\le 1+\frac{2|a|^2}{\sqrt{1+|a|^2}}
\le 1+\frac{2|a|^2}{\sqrt{1+\epsilon_i^2|a|^2}}
\qquad\forall\, a\in\mathbb R^n,
$$
and therefore
$$
\int_{\mathbb S^{n-1}} |\nabla_\tau \xi_i|\,d\mathscr H^{n-1}
+|D^s_\tau \xi_i|(\mathbb S^{n-1})
\le C(n,M).
$$
Hence, up to passing to a subsequence,
$$
\xi_i\to \xi \qquad\text{strongly in }L^p(\mathbb S^{n-1})
$$
for some $1<p=p(n)<2$, and then a.e. on $\mathbb S^{n-1}$.

For each $i$, let $c_i\in\mathbb R$ be such that
$$
\int_{\mathbb S^{n-1}} T_{\epsilon_i}(\xi_i-c_i)\,d\mathscr H^{n-1}=0.
$$
Set
$$
\tilde\xi_i:=\xi_i-c_i,\qquad
\psi_i:=T_{\epsilon_i}(\tilde\xi_i),\qquad
\eta_i:=\tilde\xi_i-\psi_i.
$$
Then the argument of Case 2 above applies verbatim to $(\tilde\xi_i,\psi_i,\eta_i)$, and gives the analogues of \eqref{two terms}, \eqref{first part}, \eqref{psi W11}, \eqref{eta controlled by xi}, \eqref{control 1}, and \eqref{eta W11}, with constants independent of $i$.

In particular, from \eqref{first part}, \eqref{psi W11}, and the $BV$ chain
rule for the truncation, we obtain
$$
\sup_i \left(
\int_{\mathbb S^{n-1}} |\psi_i|^2\,d\mathscr H^{n-1}
+
\int_{\mathbb S^{n-1}} |\psi_i\nabla_\tau \psi_i|\,d\mathscr H^{n-1}
+\epsilon_i^{-1}|D^s_\tau\psi_i|(\mathbb S^{n-1})
\right)<\infty.
$$
Moreover $\|D_\tau\psi_i\|(\mathbb S^{n-1})\le
\|D_\tau\xi_i\|(\mathbb S^{n-1})$. Hence $\psi_i$ is compact in $L^1$, while
$\psi_i^2$ is compact in $L^1$, since
$|D_\tau(\psi_i^2)|\le C|\psi_i\nabla_\tau\psi_i|\,\HH^{n-1}
+C\epsilon_i^{-1}|D^s_\tau\psi_i|$. Thus, up to extracting a further subsequence,
\begin{equation}\label{strong L2 psi}
\text{$\psi_i$ converges strongly in $L^2(\mathbb S^{n-1})$ to some function $v$.}
\end{equation}
Since $\epsilon_i\psi_i\to 0$ a.e. on $\mathbb S^{n-1}$ and
$$\frac{|\psi_i|^2}{\sqrt{1+\ez_i^2|\psi_i|^2}}\le |\psi_i|^2, $$
Vitali's convergence theorem yields
\begin{equation}\label{strong L2 psi2}
\lim_{i\to \infty}\int_{\mathbb S^{n-1}}  \frac{|\psi_i|^2}{\sqrt{1+\ez_i^2|\psi_i|^2}} \,d\mathscr H^{n-1}=  \int_{\mathbb S^{n-1}} {|v|^2}  \,d\mathscr H^{n-1}.
\end{equation}

Let
$$
E_i:=\{\eta_i\neq 0\}=\{\epsilon_i|\tilde\xi_i|>1\}\subset \{|\psi_i|=\epsilon_i^{-1}\}.
$$
Then, by \eqref{strong L2 psi},
\begin{equation}\label{measure to 0}
\epsilon_i^{-2}|E_i|
=
\int_{E_i} |\psi_i|^2\,d\mathscr H^{n-1}
\le
\int_{\{|\psi_i|\ge \epsilon_i^{-1}\}} |\psi_i|^2\,d\mathscr H^{n-1}
\longrightarrow 0.
\end{equation}
Moreover, the analogue of \eqref{eta controlled by xi} gives
$$
\epsilon_i^{-1}\int_{\mathbb S^{n-1}}|\eta_i|\,d\mathscr H^{n-1}\le C(n,M),
$$
hence
$$
\int_{\mathbb S^{n-1}}|\eta_i|\,d\mathscr H^{n-1}\to 0.
$$
Since
$$
\int_{\mathbb S^{n-1}}\xi_i\,d\mathscr H^{n-1}
=
c_i\,\mathscr H^{n-1}(\mathbb S^{n-1})+\int_{\mathbb S^{n-1}}\eta_i\,d\mathscr H^{n-1},
$$
and the averages of $\xi_i$ are uniformly bounded, it follows that $\{c_i\}$ is bounded. Up to a subsequence,
$
c_i\to c\in\mathbb R.
$
Then, by \eqref{strong L2 psi},
$$
\psi_i+c_i \to v+c \qquad\text{strongly in }L^2(\mathbb S^{n-1}).
$$
Since $\eta_i\to 0$ in $L^1(\mathbb S^{n-1})$ and $\xi_i=\psi_i+c_i+\eta_i$, while $\xi_i\to\xi$ in measure, we conclude that
$$
\xi=v+c \qquad\text{a.e. on }\mathbb S^{n-1}.
$$

We now show that
$$
\lim_{i\to\infty} \mathcal  A_{\epsilon_i}(\xi_i)=\int_{\mathbb S^{n-1}} |\xi|^2\,d\mathscr H^{n-1}.
$$
On $\mathbb S^{n-1}\setminus E_i$, one has $\eta_i=0$, hence $\xi_i=\psi_i+c_i$ there. Therefore
$$
|\mathcal A_{\epsilon_i}(\xi_i)-\mathcal A_{\epsilon_i}(\psi_i+c_i)|
\le
\int_{E_i}\frac{|\xi_i|^2}{\sqrt{1+\epsilon_i^2|\xi_i|^2}}\,d\mathscr H^{n-1}
+
\int_{E_i}\frac{|\psi_i+c_i|^2}{\sqrt{1+\epsilon_i^2|\psi_i+c_i|^2}}\,d\mathscr H^{n-1}.
$$
Since $|\psi_i|=\epsilon_i^{-1}$ on $E_i$ and $c_i$ is bounded, for $i$ large enough
$$
|\psi_i+c_i|\le 2\epsilon_i^{-1}\qquad\text{on }E_i,
$$
and therefore
$$
\int_{E_i}\frac{|\psi_i+c_i|^2}{\sqrt{1+\epsilon_i^2|\psi_i+c_i|^2}}\,d\mathscr H^{n-1}
\le C\,\epsilon_i^{-2}|E_i|\to 0
$$
by \eqref{measure to 0}. Also, on $E_i$,
$$
\xi_i=(\psi_i+c_i)+\eta_i.
$$
Since $|\psi_i+c_i|\ge \frac{1}{2}\epsilon_i^{-1}$ on $E_i$ for $i$ large, one has
$$
\frac{|\xi_i|^2}{\sqrt{1+\epsilon_i^2|\xi_i|^2}}
\le C\frac{|\psi_i+c_i|^2}{\sqrt{1+\epsilon_i^2|\psi_i+c_i|^2}}+C\epsilon_i^{-1}|\eta_i|
\qquad\text{on }E_i.
$$
Hence
$$
\int_{E_i}\frac{|\xi_i|^2}{\sqrt{1+\epsilon_i^2|\xi_i|^2}}\,d\mathscr H^{n-1}
\le
C\int_{E_i}\frac{|\psi_i+c_i|^2}{\sqrt{1+\epsilon_i^2|\psi_i+c_i|^2}}\,d\mathscr H^{n-1}
+
C\epsilon_i^{-1}\int_{\mathbb S^{n-1}}|\eta_i|\,d\mathscr H^{n-1}.
$$
If SubCase 2.1 occurs along a subsequence, then by the analogue of \eqref{eta small2} and \eqref{strong L2 psi},
$$
\epsilon_i^{-1}\int_{\mathbb S^{n-1}}|\eta_i|\,d\mathscr H^{n-1}\to 0.
$$
If SubCase 2.2.1 or SubCase 2.2.2 occurs along a subsequence, then the analogues of \eqref{control 1} or \eqref{eta W11} imply that $\epsilon_i^{-1}\eta_i$ is uniformly bounded in $BV(\mathbb S^{n-1})$; since its support is contained in $E_i$ and $|E_i|\to 0$, compactness in $L^1$ implies
$$
\epsilon_i^{-1}\eta_i\to 0 \qquad\text{strongly in }L^1(\mathbb S^{n-1}).
$$
Thus, in all cases,
$$
\epsilon_i^{-1}\int_{\mathbb S^{n-1}}|\eta_i|\,d\mathscr H^{n-1}\to 0,
$$
and consequently
$$
\mathcal A_{\epsilon_i}(\xi_i)-\mathcal A_{\epsilon_i}(\psi_i+c_i)\to 0.
$$
Finally, since $\psi_i+c_i\to \xi$ strongly in $L^2(\mathbb S^{n-1})$ and
$$
\frac{|\psi_i+c_i|^2}{\sqrt{1+\epsilon_i^2|\psi_i+c_i|^2}}
\le |\psi_i+c_i|^2,
$$
Vitali's convergence theorem again yields
$$
\lim_{i\to \infty} \mathcal A_{\epsilon_i}(\psi_i+c_i)
=
\int_{\mathbb S^{n-1}} |\xi|^2\,d\mathscr H^{n-1}.
$$
Combining the last two formulae, we obtain \eqref{eq:xi_i to xi}. This concludes the proof.
 

%Now we prove the second part of the lemma. 
%Now $\|\xi_i\|_{L^{\infty}(\mathbb S^{n-1})}\le 1$ implies 
%$$\int_{\mathbb S^{n-1}} |\xi_i-\xi_j|^2\, d\mathscr H^{n-1}\le 2^{2-p} \int_{\mathbb S^{n-1}} |\xi_i-\xi_j|^p\, d\mathscr H^{n-1},$$
%which gives that $\xi_i$ a Cauchy sequence in $L^2(\mathbb S^{n-1})$. Thus we conclude the lemma. 
%Now for the general case, we write
%$\xi= \xi_1+\xi_2$, where
%$$\xi_1=\min\left\{ \xi ,\, \frac{1}{\ez}\right\},\, \quad \xi_2=\xi-\xi_1=\left(\xi - \frac{1}{\ez}\right)_+. $$
%Since
%$$ \int_{\mathbb S^{n-1}} |\xi|^2 \,d\mathscr H^{n-1} \le 2 \left( \int_{\mathbb S^{n-1}} |\xi_1|^2\,d\mathscr H^{n-1}+\int_{\mathbb S^{n-1}} |\xi_2|^2  \,d\mathscr H^{n-1}\right)=: 2(I_1+I_2).$$
%
%Suppose first that $I_1\ge I_2$. Then  by our result in the special case
%\begin{align*}
%\int_{\mathbb S^{n-1}} |\xi|^2 \,d\mathscr H^{n-1} \le  4 I_1& \leq C(n)  \int_{\mathbb S^{n-1}} \frac{|\nabla_\tau  \xi_1|^2}{\sqrt{1+\ez^2|\nabla_\tau \xi_1|^2}} \,d\mathscr H^{n-1}\\
%& \leq  C(n)  \int_{\mathbb S^{n-1}} \frac{|\nabla_\tau  \xi|^2}{\sqrt{1+\ez^2|\nabla_\tau \xi|^2}} \,d\mathscr H^{n-1},
%\end{align*}
%where we applied the fact that $\nabla_\tau \xi_1=\nabla_\tau \xi$ almost everywhere in $\xi< \frac{1}{\ez}$ and $\nabla_\tau \xi_1=0$ when $\xi\ge \frac{1}{\ez}.$
%
%Now we consider the second case when $I_2> I_1$. We write
%$$V=\{\ez |\nabla_\tau \xi|\le 1\},\, W= \{\ez |\nabla_\tau \xi|> 1\}. $$
%
%Assume first that
%$$\int_{W} |\nabla_\tau\xi_2| \,d\mathscr H^{n-1} \le  \ez \int_{\mathbb S^{n-1}} |\xi_2|^2 \,d\mathscr H^{n-1}.$$
%Then Poincar\'e inequality tells that
%\begin{align*}
%\int_{\mathbb S^{n-1}} |\xi_2|^2 \,d\mathscr H^{n-1}& \leq  \int_{\mathbb S^{n-1}} |D\xi_2|^2 \,d\mathscr H^{n-1}\\
%& \leq  \int_{V\cap \{\ez\xi>1\}} |D\xi_2|^2 \,d\mathscr H^{n-1} + \int_{W\cap \{\ez\xi>1\}} |D\xi_2|^2 \,d\mathscr H^{n-1}\\
%& \leq  \int_{V \cap \{\ez\xi>1\}} |D\xi_2|^2 \,d\mathscr H^{n-1} + 
%\end{align*}
%and then
%$$\int_{\{\ez \xi>1 \}\cap V}|\nabla_\tau \xi_2|^2 \,d\mathscr H^{n-1}\le 2 \int_{\{\ez \xi>1 \}\cap V}\frac{|\nabla_\tau \xi_2|^2}{\sqrt{1+\ez^2|\nabla_\tau \xi_2|^2 }} \,d\mathscr H^{n-1}.$$
%On the other hand,   
%\begin{align*}
% & \leq  C(n)  \int_{\{\ez \xi>\zeta \}\cap W}\frac{|\nabla_\tau \xi_2|}{ \ez  } \,d\mathscr H^{n-1}\\
%& \leq  C(n) \int_{\{\ez \xi>\zeta \}\cap W}\frac{|\nabla_\tau \xi_2|^2}{\sqrt{1+\ez^2|\nabla_\tau \xi_2|^2 }} \,d\mathscr H^{n-1}
%\end{align*}
\end{proof}

\end{document}
